%==============================================================================
%  Plan de projet expérimental — IFT-7026
%  Université Laval
%
%  Sujet : ordonnancement Linux assisté par apprentissage (candidate-based),
%          de l'émulateur au noyau réel via sched_ext.
%
%  Compilation : pdflatex (2 passes) ou latexmk -pdf plan_projet_IFT-7026.tex
%==============================================================================
\documentclass[11pt,a4paper]{article}

\usepackage[utf8]{inputenc}
\usepackage[T1]{fontenc}
\usepackage[french]{babel}
\usepackage{lmodern}
\usepackage[margin=2.5cm]{geometry}
\usepackage{enumitem}
\usepackage{booktabs}
\usepackage{array}
\usepackage{tabularx}
\usepackage{xcolor}
\usepackage{titlesec}
\usepackage[hidelinks]{hyperref}

% --- Mise en forme des titres de section -------------------------------------
\definecolor{ulgris}{RGB}{60,60,60}
\titleformat{\section}{\large\bfseries\color{ulgris}}{\thesection.}{0.5em}{}
\titlespacing*{\section}{0pt}{1.4em}{0.6em}

% --- Métadonnées du document -------------------------------------------------
\newcommand{\projetTitre}{Plan de projet expérimental — IFT-7026}
\newcommand{\projetSujet}{Un ordonnanceur Linux appris par sélection de
  candidats : de l'émulateur au noyau réel via \texttt{sched\_ext}}
\newcommand{\projetSuperviseur}{Prof.\ Mohamed Aymen Saied}
\newcommand{\projetDirecteur}{Prof.\ Pascal Tesson}
\newcommand{\projetEtudiantUn}{Léopold Chappuis}
\newcommand{\projetEtudiantDeux}{Alexandre Eberhardt}
\newcommand{\projetSession}{Session d'automne 2026}
\newcommand{\projetUniversite}{Université Laval}

\begin{document}

%==============================================================================
%  En-tête
%==============================================================================
\begin{center}
  {\LARGE\bfseries \projetTitre}\\[0.4em]
  {\large \projetUniversite}
\end{center}

\vspace{0.8em}
\noindent
\begin{tabular}{@{}ll@{}}
  \textbf{Étudiants}              & \projetEtudiantUn{}, \projetEtudiantDeux \\
  \textbf{Superviseur}            & \projetSuperviseur \\
  \textbf{Directeur de programme} & \projetDirecteur \\
  \textbf{Session}                & \projetSession \\
\end{tabular}

\vspace{0.6em}
\noindent\textbf{Sujet :} \projetSujet

\vspace{0.4em}
\hrule
\vspace{1.2em}

%==============================================================================
\section{Objectif du projet}
%==============================================================================
Ce projet cherche à savoir si une politique d'ordonnancement \emph{apprise} peut
faire mieux que simplement imiter les ordonnanceurs classiques du noyau Linux
(CFS, EEVDF), tout en restant assez légère pour être \textbf{déployée dans un
noyau réel}. Le critère principal est l'amélioration du temps de réponse,
\emph{sous contraintes} : sans famine et sans dégradation notable des autres
métriques (équité, commutations de contexte, turnaround) — optimiser le temps
de réponse seul serait trivial au prix de la famine des tâches longues. Le
projet prolonge des résultats préliminaires obtenus dans le
cadre du cours GLO-7030, où une formulation \emph{candidate-based} évaluée en
boucle fermée a déjà montré une réduction d'environ \mbox{$10{,}5\,\%$} du
temps de réponse moyen face à CFS dans un émulateur.

\medskip
Les objectifs spécifiques sont :
\begin{itemize}[leftmargin=1.5em,itemsep=3pt]
  \item Améliorer la qualité du signal d'apprentissage : aujourd'hui le modèle
    apprend à imiter un \emph{teacher} fait de règles fixes. Il sera remplacé
    par un \textbf{teacher quasi-optimal} qui, pour chaque décision, essaie
    plusieurs choix dans l'émulateur et regarde quelques pas en avant pour
    identifier le meilleur. Le modèle apprend ensuite à reproduire ces bonnes
    décisions sans avoir accès au futur (il n'utilise que l'information
    disponible au moment de décider).
  \item Corriger le décalage entre entraînement et utilisation
    (\emph{covariate shift}) : le modèle est entraîné sur des situations
    produites par CFS, mais une fois aux commandes il crée ses propres
    situations, jamais vues à l'entraînement. Le remède (\textbf{réentraînement
    itératif de type DAgger}) consiste à laisser le modèle ordonnancer, à faire
    étiqueter par le \emph{teacher} les situations ainsi créées, puis à
    réentraîner sur ces nouvelles données.
  \item Caractériser le compromis \textbf{précision / latence d'inférence /
    nombre de paramètres} par distillation (front de Pareto), la latence étant
    le principal facteur limitant côté systèmes.
  \item \textbf{Déployer} le modèle dans un vrai noyau Linux, pas seulement
    dans l'émulateur. Le mécanisme \texttt{sched\_ext} (Linux $\geq$ 6.12)
    permet de brancher un ordonnanceur personnalisé sans recompiler le noyau :
    sur le modèle de \texttt{scx\_rustland}, un petit composant dans le noyau
    transmet chaque décision à un programme en espace utilisateur, qui
    interrogera ici le modèle distillé pour choisir la prochaine tâche. On
    mesurera alors, sur de vraies charges (p.~ex.\ compilation, serveur sous
    requêtes), si la machine répond mieux qu'avec l'ordonnanceur par défaut
    (EEVDF), avec les mêmes métriques que dans l'émulateur.
\end{itemize}

%==============================================================================
\section{Résumé}
%==============================================================================
L'ordonnanceur du noyau Linux prend des milliers de décisions par seconde, qui
déterminent latence, équité et utilisation du CPU. Ces décisions reposent sur des
heuristiques affinées sur des décennies (CFS, puis EEVDF depuis Linux~6.6), qui
n'optimisent pas explicitement des objectifs comme le temps de réponse. La
question posée est de savoir si un modèle appris peut reconstruire, voire
améliorer, ces décisions à partir de l'état d'exécution. L'application d'un
modèle appris à la décision d'ordonnancement, et plus encore son déploiement
dans un noyau réel, reste peu explorée dans la
littérature : les travaux existants (dont \textit{KernelOracle}~[3])
[1--4] se limitent à la prédiction hors ligne, sans déploiement ni
évaluation en boucle fermée dans un noyau réel.

\medskip
Le projet s'appuie sur une base existante (GLO-7030) comportant deux parties :
(i) un pipeline d'imitation séquentielle (modèle récurrent LSTM) sur traces
\texttt{sched\_switch} réelles collectées avec \texttt{perf}, conservé ici comme
\textbf{baseline} (\og imitation pure \fg{}, proche de KernelOracle) ; et (ii) une
formulation \emph{candidate-based} apprenant à choisir une tâche parmi des
candidats à partir d'un état d'ordonnancement riche, via un encodeur à
attention sur l'ensemble des tâches et l'historique de décisions, évaluée en
boucle fermée dans un émulateur. C'est cette seconde partie qui constitue la contribution
principale.

\medskip
Le projet proposé vise à corriger les limites de ces travaux préliminaires : un
\emph{teacher} heuristique trop simple, un \emph{covariate shift} non traité, et une
évaluation restreinte à l'émulateur. Il propose un \emph{teacher} quasi-optimal,
un réentraînement DAgger, une analyse du compromis précision/latence par
distillation, puis un \textbf{déploiement dans un noyau réel via
\texttt{sched\_ext}}. L'accent est mis autant sur la \textbf{méthodologie
d'évaluation} (protocole en boucle fermée, baselines EEVDF/CFS, métriques dont la
latence de queue p99, tests statistiques multi-graines) et la
\textbf{reproductibilité} du pipeline que sur l'apprentissage lui-même :
il ne s'agit pas seulement d'entraîner un modèle, mais de mesurer
précisément et de façon reproductible ce qu'il apporte. Les entraînements
s'appuieront sur les ressources de calcul
(serveur ICE, via le compte de recherche du superviseur) ; le déploiement
\texttt{sched\_ext} ciblera une machine disposant d'un noyau Linux~6.12
(\texttt{CONFIG\_SCHED\_CLASS\_EXT}) avec accès privilégié. Le résultat final
visé est un rapport au format article, soumissible à un atelier de type
\emph{ML for Systems} (repli : EuroMLSys, APSys).

%==============================================================================
\section{Contenu et travaux à effectuer}
%==============================================================================
Le projet est structuré en trois phases.

\medskip
\noindent\textbf{Phase 0 — Consolidation et protocole expérimental}
\begin{itemize}[leftmargin=1.5em,itemsep=2pt]
  \item Nettoyage et reproductibilité du code existant (pipeline, émulateur).
  \item Unification du \emph{baseline} de comparaison sur EEVDF (cohérence avec
    le noyau réel ; correction de l'écart CFS/EEVDF des travaux préliminaires).
  \item Définition d'une suite de charges de travail, des métriques (temps de
    réponse moyen et \textbf{latence de queue p99}, turnaround, commutations de
    contexte, équité) et d'un protocole statistique (multi-graines,
    significativité).
  \item Revue de littérature et positionnement.
\end{itemize}

\noindent\textbf{Phase 1 — Contribution méthodologique (émulateur)}
\begin{itemize}[leftmargin=1.5em,itemsep=2pt]
  \item \emph{Teacher} quasi-optimal à lookahead et distillation vers un modèle
    causal.
  \item Réentraînement itératif DAgger sur les trajectoires auto-induites.
  \item Front de Pareto précision / latence d'inférence / nombre de paramètres.
  \item Évaluation en boucle fermée vs EEVDF et CFS, ablations de l'architecture
    \emph{candidate-based}, généralisation à des charges non vues.
\end{itemize}

\noindent\textbf{Phase 2 — Déploiement noyau réel}
\begin{itemize}[leftmargin=1.5em,itemsep=2pt]
  \item Intégration d'un ordonnanceur \texttt{sched\_ext} en espace utilisateur
    (type \texttt{scx\_rustland}) interrogeant le modèle distillé.
  \item Mesures en boucle fermée sur charges réelles vs EEVDF ; analyse du budget
    de latence d'inférence.
  \item \emph{Repli} : si le déploiement \texttt{sched\_ext} s'avère bloquant,
    bascule vers un ordonnanceur \og orientable \fg{} (politique unique
    conditionnée par un vecteur de poids d'objectifs), évalué en émulateur.
\end{itemize}

\medskip
\noindent Liste détaillée des travaux :
\begin{enumerate}[leftmargin=1.6em,itemsep=3pt]
  \item Effectuer la revue de littérature (imitation, prédiction séquentielle,
    RL/imitation pour l'ordonnancement, \texttt{sched\_ext}) et positionner la
    contribution.
  \item Consolider le code et fixer le protocole expérimental (charges,
    métriques dont p99, statistiques) ; aligner le \emph{baseline} sur EEVDF.
  \item Concevoir le \emph{teacher} quasi-optimal à lookahead et le comparer au
    \emph{teacher} heuristique des travaux préliminaires.
  \item Entraîner par distillation et tracer le front de Pareto précision /
    latence / paramètres ; intégrer le réentraînement DAgger.
  \item Évaluer en boucle fermée dans l'émulateur (vs EEVDF/CFS), réaliser les
    ablations et les tests de généralisation ; rédiger le rapport de mi-session.
  \item Intégrer la politique distillée dans un ordonnanceur \texttt{sched\_ext}
    et mesurer sur un noyau réel face à EEVDF ; analyser le budget de latence.
  \item Rédiger le rapport final au format article et identifier les pistes de
    recherche futures.
\end{enumerate}

\medskip
\noindent\textbf{Faisabilité et portée.} La fondation déjà en place (code,
émulateur et premiers résultats issus de GLO-7030) réduit le risque et concentre
l'effort sur la contribution. Les phases~0 et~1, réalisables dans l'enveloppe de
\mbox{5--6\,h/semaine} par personne sur la session, constituent le \textbf{livrable
garanti} du projet. Le déploiement dans un noyau réel (phase~2) en est
l'\textbf{objectif ambitieux} : en cas de blocage technique sur
\texttt{sched\_ext}, le repli « orientable » décrit ci-dessus assure une
contribution expérimentale complète et évaluable en émulateur.

%==============================================================================
\section{Modalités d'évaluation}
%==============================================================================
\noindent
\renewcommand{\arraystretch}{1.3}
\begin{tabularx}{\linewidth}{@{}Xr@{}}
  \toprule
  \textbf{Composante} & \textbf{Pondération} \\
  \midrule
  Présentation orale — état de l'art et positionnement & 15 \% \\
  \addlinespace
  Production expérimentale et reproductibilité (Phase 1 : code, émulateur, \emph{benchmark}) & 15 \% \\
  \addlinespace
  Rapport de mi-session (présenté à la présentation intermédiaire) & \textbf{35 \%} \\
  \quad Description de la problématique et de la méthodologie & 10 \% \\
  \quad Analyse quantitative (boucle fermée, ablations, Pareto) & 15 \% \\
  \quad Analyse qualitative (analyse d'erreurs, limites) & 10 \% \\
  \addlinespace
  Rapport final (format article, soutenu à la présentation finale) & \textbf{35 \%} \\
  \quad Nouveaux aspects étudiés depuis la mi-session & 10 \% \\
  \quad Analyse quantitative (\texttt{sched\_ext} vs EEVDF, latence) & 15 \% \\
  \quad Analyse qualitative & 10 \% \\
  \midrule
  \textbf{Total} & \textbf{100 \%} \\
  \bottomrule
\end{tabularx}

\medskip
\noindent\small Les présentations intermédiaire et finale ne sont pas pondérées
séparément : elles constituent la soutenance orale des rapports correspondants
et sont évaluées au sein de ceux-ci. La présentation orale de l'état de l'art
(début de session) est la seule présentation pondérée de façon autonome. Les
présentations peuvent se tenir \textbf{en visioconférence}.

%==============================================================================
\section{Échéancier}
%==============================================================================
\noindent
\renewcommand{\arraystretch}{1.3}
\begin{tabularx}{\linewidth}{@{}Xl@{}}
  \toprule
  \textbf{Étape} & \textbf{Date cible} \\
  \midrule
  Début du projet (Phase 0)                              & Début septembre 2026 \\
  Présentation orale — état de l'art et positionnement   & Fin septembre 2026 \\
  Livrable 1 — résultats Phase 1 + rapport de mi-session & Fin octobre 2026 \\
  Présentation intermédiaire (soutenance du rapport de mi-session) & Mi-novembre 2026 \\
  Livrable 2 — déploiement \texttt{sched\_ext} (Phase 2) & Début décembre 2026 \\
  Rapport final (format article) + présentation finale   & Mi-décembre 2026 \\
  \bottomrule
\end{tabularx}

%==============================================================================
\section{Livrables}
%==============================================================================
\begin{itemize}[leftmargin=1.5em,itemsep=2pt]
  \item Code reproductible : pipeline d'apprentissage, émulateur, intégration
    \texttt{sched\_ext}.
  \item Suite de charges de travail / \emph{benchmark} et jeux de décisions
    candidates utilisés pour les expérimentations.
  \item Rapport de mi-session et rapport final.
  \item Présentations : présentation orale (état de l'art), présentation
    intermédiaire et présentation finale.
  \item Draft d'article soumissible (cible : \emph{ML for Systems} ;
    repli : EuroMLSys, APSys).
\end{itemize}

%==============================================================================
\section{Références (préliminaires)}
%==============================================================================
\begin{enumerate}[leftmargin=2em,itemsep=4pt,label={[\arabic*]}]
  \item Jingde Chen, Subho S.\ Banerjee, Zbigniew T.\ Kalbarczyk et Ravishankar K.\ Iyer.
    \emph{Machine learning for load balancing in the linux kernel}. Dans
    \textit{Proceedings of the 11th ACM SIGOPS Asia-Pacific Workshop on Systems},
    APSys~'20, p.~67--74, New York, NY, USA, 2020. Association for Computing Machinery.
  \item Anusha G.~H., Minal K., B.~P.\ Varsha, Geethishree Mishra, Meghana G.~K.\ et
    Madhusudhan K.~N. \emph{ML engine for linux scheduler}. \textit{Quest Journals,
    Journal of Software Engineering and Simulation}, vol.~9, n\textsuperscript{o}~7,
    p.~20--29, 2023.
  \item Sampanna Yashwant Kahu. \emph{KernelOracle: Predicting the linux scheduler's
    next move with deep learning}, 2025.
  \item Atul Negi et P.\ Kishore Kumar. \emph{Applying machine learning techniques to
    improve linux process scheduling}. Dans \textit{TENCON 2005 -- 2005 IEEE Region 10
    Conference}, p.~1--6, 2005.
  \item Stéphane Ross, Geoffrey J.\ Gordon et Drew Bagnell. \emph{A reduction of
    imitation learning and structured prediction to no-regret online learning}. Dans
    \textit{Proceedings of the Fourteenth International Conference on Artificial
    Intelligence and Statistics (AISTATS)}, vol.~15 de \textit{Proceedings of Machine
    Learning Research}, p.~627--635, 2011.
  \item Vasuki Shankar. \emph{Machine learning for linux kernel optimization: Current
    trends and future directions}. \textit{International Journal of Computer Sciences
    and Engineering}, vol.~13, p.~56--64, 03~2025.
  \item Ion Stoica et Hussein Abdel-Wahab. \emph{Earliest Eligible Virtual Deadline
    First: A flexible and accurate mechanism for proportional share resource
    allocation}. Rapport technique TR-95-22, Old Dominion University, 1995.
  \item \emph{Sched\_ext: BPF extensible scheduler class}. Documentation du noyau
    Linux. \url{https://docs.kernel.org/scheduler/sched-ext.html} (consulté en 2026).
\end{enumerate}

\end{document}
